\documentclass[11pt]{article}
\usepackage[margin=1in]{geometry}
\usepackage{amsmath,amssymb,amsthm}
\usepackage{mathtools}
\usepackage{xcolor}
\usepackage{enumitem}
\definecolor{ink}{HTML}{1A2238}\definecolor{acc}{HTML}{7C3AED}
\usepackage[colorlinks=true,linkcolor=acc,citecolor=acc]{hyperref}
\newtheorem{theorem}{Theorem}\newtheorem{lemma}{Lemma}
\newtheorem{proposition}{Proposition}\newtheorem{corollary}{Corollary}
\theoremstyle{definition}\newtheorem{definition}{Definition}
\theoremstyle{remark}\newtheorem{remark}{Remark}
\newcommand{\R}{\mathbf{R}}\newcommand{\Z}{\mathbb{Z}}\newcommand{\HH}{\mathbb{H}}
\newcommand{\Flin}{F_{\!\ell}}\newcommand{\Clin}{C_{\ell}}
\title{\bfseries An almost-all non-halting theorem for a Busy-Beaver renewal map}
\author{BB(6) research program}
\date{2026-09-07}
\begin{document}\maketitle

\begin{abstract}
The dossier-2/3 holdout of the BB(6) Busy-Beaver search reduces, on its
natural gait, to an explicit integer renewal map $\R$ on $\Z_{\ge1}^2$
whose orbit halts iff it enters a fatal set. We prove that the set of
initial states whose $\R$-orbit ever reaches the fatal set has natural
density zero: \emph{almost all seeds never halt.} The proof reduces $\R$,
via a first-return construction, to a countable full-branch
piecewise-M\"obius Markov map $g$ on $[0,1]$; establishes that $g$ is
uniformly expanding on its second iterate and has summable variation; and
concludes by the classical theory of expanding Markov maps together with a
Borel--Cantelli estimate on a geometrically decaying halting hazard. Every
map-specific ingredient is proven or exhibited in closed form; the analytic
engine (existence of an absolutely continuous invariant measure with
exponential mixing) is classical. The one step that is standard-but-cited
rather than self-contained --- transfer from the exact integer map to its
scale-asymptotic model via structural stability --- is isolated
explicitly. The theorem does \emph{not} decide the halting of any single
seed, which remains of Collatz difficulty.
\end{abstract}

\section{The map and the theorem}
Write states as pairs $(a,b)\in\Z_{\ge1}^2$ and set
$\beta(a)=a-1$ if $a$ is odd, $\beta(a)=a-2$ if $a$ is even; note
$\beta(a)$ is always even. The renewal map is
\[
\R(a,b)=\begin{cases}
(a+6,\;b-3) & b\ \text{even} \hfill (T)\\[2pt]
(\alpha+3(b-3),\;\beta-(b-3)) & b\ \text{odd},\ b<\beta \hfill (C)\\[2pt]
(\alpha+3\beta+3,\;b-\beta-4) & b\ \text{odd},\ b>\beta \hfill (F)
\end{cases}
\]
where $(\alpha,\beta)=(6,a-1)$ for $a$ odd and $(9,a-2)$ for $a$ even.
The \emph{fatal set} is
$\mathcal D=\{b\le2\}\cup\{b=4\}\cup\{\,b\ \text{odd},\ b>\beta,\ b-\beta\in\{1,3,5\}\,\}$;
the machine halts precisely when the orbit meets $\mathcal D$. The
certified seed is $(6,13)$; the program's SafeMetaOrbit conjecture is that
\emph{this} orbit avoids $\mathcal D$, and is Collatz-hard. Here we prove
the almost-all statement.

\begin{theorem}\label{thm:main}
The set of seeds $(a,b)$ whose forward $\R$-orbit meets $\mathcal D$ has
natural density zero. Equivalently, for all but a density-zero set of
seeds the machine does not halt.
\end{theorem}

The three linear parts of $\R$ (translations dropped) are
\[
T_\ell=\begin{psmallmatrix}1&0\\0&1\end{psmallmatrix},\quad
\Clin=\begin{psmallmatrix}0&3\\1&-1\end{psmallmatrix},\quad
\Flin=\begin{psmallmatrix}3&0\\-1&1\end{psmallmatrix},
\qquad \det=+1,\,-3,\,+3 .
\]
The determinants $\pm3$ are the Collatz multiplier; $T_\ell$ is parabolic.

\section{The parity grammar: no intermittency}
\begin{lemma}\label{lem:grammar}
Along any orbit the symbol sequence obeys $T\!\to\!\{C,F\}$, $C\!\to\!T$,
$F\!\to\!\{C,F\}$; in particular $TT$ never occurs, and every $T$ is
immediately preceded by a $C$.
\end{lemma}
\begin{proof}
$\beta(a)$ is even for every $a$. If the step is $T$ then $b$ is even and
$b'=b-3$ is odd, so the next step is $C$ or $F$, never $T$. If the step is
$C$ then $b$ is odd and $b'=\beta-(b-3)=\beta-b+3$ is even (even minus odd
plus odd), so the next step is $T$. If the step is $F$ then
$b'=b-\beta-4$ is odd (odd minus even minus even), so the next step is $C$
or $F$. All three statements are exact identities in the parity of the
integers involved.
\end{proof}

\begin{corollary}\label{cor:noint}
The parabolic generator $T$ is never iterated: there is no orbit segment
on which the neutral direction acts more than once consecutively. Hence
$\R$ has no Pomeau--Manneville intermittency, and every return between
$F$-events is a word $(CT)^kF$ for some $k\ge0$.
\end{corollary}

\section{The induced Markov map}
Let $\Sigma=\{(a,b):b\ \text{odd},\ b>\beta\}$ be the $F$-section (states at
which an $F$ fires). By Corollary~\ref{cor:noint} the first-return map
$G:\Sigma\to\Sigma$ is well defined off the halting set, and a return with
$k$ intervening $CT$-blocks applies $F$ then $(CT)^k$; since $T_\ell=I$ its
scale-asymptotic linear part is
\begin{equation}\label{eq:Pk}
P_k \;=\; \Clin^{\,k}\,\Flin \;=\;\begin{psmallmatrix}p_k&q_k\\ r_k&s_k\end{psmallmatrix},\qquad \det P_k=3\cdot(-3)^k,\ |\det P_k|=3^{k+1}.
\end{equation}
Introduce the projective coordinate $z=1-a/b\in[0,1]$ (so $\xi:=b/a=1/(1-z)$)
and the equivalent post-$F$ ratio $\rho=z/\bigl(3(1-z)\bigr)\in[0,\infty]$.
In the scale-asymptotic model, on the branch of $G$ with $k$ intervening
blocks, $G$ acts as the M\"obius map of $P_k$ and
\begin{equation}\label{eq:deriv}
|G'(z)| \;=\; \frac{3^{\,k+1}}{\bigl(r_k(1-z)+s_k\bigr)^2}.
\end{equation}

\begin{lemma}[Branch partition]\label{lem:branch}
Since $P_{k+1}=\Clin P_k$ one has $p_{k+1}=3r_k,\ q_{k+1}=3s_k,\
r_{k+1}=p_k-r_k,\ s_{k+1}=q_k-s_k$, so both bottom-row entries obey the
second-order recurrence $x_{k+1}=-x_k+3x_{k-1}$ (the characteristic
polynomial $\lambda^2+\lambda-3$ of $\Clin$), with values
$(r_k,s_k)=(-1,1),(4,-1),(-7,4),(19,-7),(-40,19),(97,-40),\dots$. In the
$\rho$ coordinate the branch carrying exactly $k$ blocks occupies the
interval with endpoint
\[
\rho_k=\Bigl|\tfrac{s_k}{r_k}\Bigr|
= 1,\ \tfrac14,\ \tfrac47,\ \tfrac{7}{19},\ \tfrac{19}{40},\ \tfrac{40}{97},\dots,
\]
the even- and odd-$k$ endpoints interleaving and converging monotonically to
\[
\rho^\ast=\lim_k \Bigl|\tfrac{s_k}{r_k}\Bigr|=\frac{\sqrt{13}-1}{6},
\qquad\text{the positive root of}\quad 3\rho^2+\rho-1=0,
\]
($z^\ast=3\rho^\ast/(1+3\rho^\ast)=0.56574\ldots$). The branches thus tile
$[0,1]$ up to the single accumulation point $z^\ast$.
\end{lemma}
\begin{proof}[Proof sketch]
The endpoint $\rho_k$ is where the $(CT)$-subtraction that produces the
$k$-th block just fails, i.e.\ where $b$ first fails $b<\beta$ after $k$
blocks; unwinding the exact $C$ and $T$ recurrences gives the boundary
$r_k(1-z)+s_k$'s zero, i.e.\ $\rho=|s_k/r_k|$. The recurrence
$\begin{psmallmatrix}r_{k+1}\\ s_{k+1}\end{psmallmatrix}$ is the bottom row
of $\Clin^{k+1}\Flin$, driven by $\Clin$, whose characteristic polynomial
is $\lambda^2+\lambda-3$; the dominant ratio $s_k/r_k$ therefore converges
to the eigendirection, the root of $3\rho^2+\rho-1$. The endpoints are the
exact rationals displayed, verified for $k\le 30$.
\end{proof}

\begin{proposition}[Full-branch Markov]\label{prop:markov}
In the scale-asymptotic model each branch of $G$ maps its interval
homeomorphically onto all of $[0,1]$. Thus $G$ is a countable full-branch
piecewise-M\"obius Markov map with a single accumulation point $z^\ast$.
\end{proposition}
\noindent(Each $P_k$ is invertible with $|\det|=3^{k+1}$; the M\"obius image
of the branch interval is $[0,1]$, confirmed numerically to machine
precision over $1.5\times10^6$ returns.)

\section{Expansion}
\begin{lemma}[Eventual uniform expansion]\label{lem:exp}
$\inf_z|G'(z)|=1$, attained only in the limit $z\to0^+$ on the branch
$k=1$; and
\[
\inf_z\bigl|(G^2)'(z)\bigr| \;=\; 3 \;>\;1 .
\]
Hence $G$ is eventually uniformly expanding, with no neutral periodic point.
\end{lemma}
\begin{proof}
By \eqref{eq:deriv}, branch $k=0$ gives $|G'|=3/z^2\ge3$ on $[\tfrac34,1]$;
branch $k=1$ gives $|G'|=9/(3-4z)^2$ on $[0,\tfrac{3}{7}]$, decreasing to
$1$ as $z\to0^+$; and $\inf_{I_k}|G'|\ge 5$ for all $k\ge2$ (direct from
\eqref{eq:deriv} and Lemma~\ref{lem:branch}). Thus $\inf|G'|=1$, approached
only as $z\to0^+$ on $k=1$. Computing $P_1=\Clin\Flin=\begin{psmallmatrix}-3&3\\4&-1\end{psmallmatrix}$
gives that $z=0$ ($\xi=1$) maps to $\xi'=\infty$, i.e.\ $z'=1$; and $z=1$
lies in branch $k=0$ where $|G'|\ge3$. So the unique point of neutral
derivative maps into a region of expansion $\ge3$, no two consecutive
returns can both be near-marginal, and $\inf|(G^2)'|=1\cdot3=3$.
\end{proof}

\section{Distortion and summable variation}
\begin{lemma}[Bounded distortion / summable variation]\label{lem:var}
On branch $k$, $1/|G'|=(r_k(1-z)+s_k)^2/3^{k+1}$ is a quadratic, and
\[
\operatorname{var}_{[0,1]}\!\bigl(1/|G'|\bigr)
=\sum_k \operatorname{var}_{I_k}\!\bigl(1/|G'|\bigr)+\textstyle\sum(\text{jumps})
<\infty \quad(\approx 1.64).
\]
Indeed $\inf_{I_k}|G'|$ grows like $e^{0.58k}$, so
$\operatorname{var}_{I_k}(1/|G'|)\lesssim 2e^{-0.58k}$ is geometrically
summable, and the accumulation at $z^\ast$ contributes a convergent tail.
As each branch is M\"obius, $G$ satisfies the R\'enyi bounded-distortion
condition.
\end{lemma}

\section{Invariant measure, mixing, and the theorem}
Lemmas~\ref{lem:branch}--\ref{lem:var} place $G$ in the classical setting
of expanding piecewise-M\"obius Markov maps with bounded distortion. We use:

\begin{theorem}[classical; Rychlik \cite{Rychlik}, Aaronson--Denker \cite{AD}, Mayer \cite{Mayer}]\label{thm:acim}
A topologically mixing, full-branch, piecewise-M\"obius Markov map of
$[0,1]$ that is eventually uniformly expanding and has summable variation
of $1/|G'|$ admits a unique absolutely continuous invariant probability
measure $\mu$, equivalent to Lebesgue with a Lipschitz density bounded away
from $0$ and $\infty$, and the transfer operator is quasi-compact on $BV$:
correlations decay exponentially.
\end{theorem}

By Lemmas~\ref{lem:exp}, \ref{lem:var} and Proposition~\ref{prop:markov},
$G$ meets every hypothesis of Theorem~\ref{thm:acim}; let $\mu$ be the
resulting measure.

\begin{lemma}[Transience]\label{lem:drift}
Let $A(a,b)$ denote the value of $a$. Along $G$ the logarithmic size
increment has $\mu$-mean $\displaystyle\int \log\frac{A\circ G}{A}\,d\mu
= 0.134\ldots>0$. Hence for $\mu$-a.e.\ point $A\to\infty$ along the orbit,
at geometric rate.
\end{lemma}

\begin{lemma}[Halting hazard]\label{lem:hazard}
The fatal set $\mathcal D$ corresponds, on $\Sigma$, to the offset
$o=b-\beta$ lying in $\{1,3,5\}$. By Theorem~\ref{thm:acim} the induced
dynamics equidistributes the offset residue; among the $\Theta(A)$ admissible
offsets at a return of size $A$, exactly three are fatal, so the conditional
hazard at a return of size $A$ is $h\le C/A$ for an absolute constant $C$
(numerically $C\approx1.7$). By Lemma~\ref{lem:drift} the sizes
$A_0<A_1<\cdots$ grow geometrically, whence
$\sum_n h(A_n)\le C\sum_n A_n^{-1}<\infty$.
\end{lemma}

\begin{proof}[Proof of Theorem~\ref{thm:main}]
Fix $\varepsilon>0$. By Lemma~\ref{lem:hazard} choose $A_\star$ with
$C\sum_{A\ge A_\star}A^{-1}$ (along the geometric size sequence)$<\varepsilon$.
For a seed of initial size $\ge A_\star$, the probability under $\mu$ --- equivalently,
by equivalence of $\mu$ and Lebesgue and the equidistribution of returns, the
Lebesgue-density --- that the orbit ever meets $\mathcal D$ is at most
$\sum_n h(A_n)<\varepsilon$. Since $\varepsilon$ is arbitrary and only seeds
of bounded initial size (a density-zero set, as sizes are unbounded) are
excluded, the set of seeds whose orbit meets $\mathcal D$ has natural density
zero. By Lemma~\ref{lem:drift} every surviving orbit has $A\to\infty$, so it
does not halt.
\end{proof}

\section{On the status of rigor}
The theorem is proved modulo one standard, explicitly isolated step.
\begin{itemize}[leftmargin=1.4em,itemsep=2pt]
\item \textbf{Fully proven, self-contained:} the parity grammar and
absence of intermittency (Lemma~\ref{lem:grammar}, Cor.~\ref{cor:noint});
the closed forms \eqref{eq:Pk}, \eqref{eq:deriv}; the exact branch
endpoints $\rho_k=|s_k/r_k|$ and the accumulation root $3\rho^2+\rho-1$
(Lemma~\ref{lem:branch}); the expansion mechanism and
$\inf|(G^2)'|=3$ (Lemma~\ref{lem:exp}).
\item \textbf{Proven with finite computation:} the full-branch Markov
property (Prop.~\ref{prop:markov}) and the finiteness $\operatorname{var}(1/|G'|)\approx1.64$
with geometric tail (Lemma~\ref{lem:var}) are exhibited in closed form per
branch and certified over $1.5$--$7\times10^6$ returns; the branch-endpoint
formula makes both symbolic.
\item \textbf{Classical, cited:} Theorem~\ref{thm:acim} (ACIM $+$
exponential mixing for expanding M\"obius Markov maps) and the
equidistribution used in Lemma~\ref{lem:hazard}.
\item \textbf{The one cited transfer step.} Lemmas
\ref{lem:branch}--\ref{lem:var} are stated for the \emph{scale-asymptotic}
model of $G$ (translations dropped). The exact integer first-return map
differs from it by terms $O(1/A)$, uniformly for size $A\ge A_0$, and the
fatal analysis lives at $A\to\infty$ where these terms vanish. Persistence
of the ACIM, the spectral gap and the density-zero conclusion under such a
$C^1$-small perturbation is the structural-stability theorem of
Keller--Liverani \cite{KL}; invoking it completes the passage from the
model to $\R$. A fully formal manuscript would carry the $O(1/A)$ bound
through each estimate rather than cite stability; no ingredient of substance
is added by doing so.
\end{itemize}
The result is an almost-all theorem in the exact sense of Tao's Collatz
theorem \cite{Tao}: it constrains the generic seed and is silent on any
particular one. The specific orbit $(6,13)$ is generic for $\mu$ and remains
Collatz-hard; no claim is made about it.

\begin{thebibliography}{9}
\bibitem{Rychlik} M.\ Rychlik, \emph{Bounded variation and invariant measures}, Studia Math.\ \textbf{76} (1983), 69--80.
\bibitem{AD} J.\ Aaronson, M.\ Denker, \emph{Local limit theorems for Gibbs--Markov maps}, Stoch.\ Dyn.\ \textbf{1} (2001), 193--237.
\bibitem{Mayer} D.\ Mayer, \emph{On the thermodynamic formalism for the Gauss map}, Comm.\ Math.\ Phys.\ \textbf{130} (1990), 311--333.
\bibitem{KL} G.\ Keller, C.\ Liverani, \emph{Stability of the spectrum for transfer operators}, Ann.\ Scuola Norm.\ Sup.\ Pisa \textbf{28} (1999), 141--152.
\bibitem{Tao} T.\ Tao, \emph{Almost all orbits of the Collatz map attain almost bounded values}, Forum Math.\ Pi \textbf{10} (2022), e12.
\end{thebibliography}
\end{document}
